\documentclass{article}
\usepackage[utf8]{inputenc}
\usepackage[T1]{fontenc}
\usepackage{graphicx}
\usepackage{algorithm}
\usepackage{algpseudocode}

\usepackage{listings}
\usepackage{xcolor}
\usepackage[margin=2.5cm]{geometry}
\usepackage{float}

\usepackage{amssymb} %Ajoute des symboles mathématiques comme R, Z, forall, exists, etc....

\lstdefinestyle{monstyle}{
    backgroundcolor=\color{white},
    commentstyle=\color{gray},
    keywordstyle=\color{blue},
    numberstyle=\tiny\color{gray},
    basicstyle=\ttfamily\small,
    breaklines=true,
    numbers=left,
    numbersep=5pt,
    frame=single,
}
\lstset{style=monstyle}


\title{Rapport TD 1 : Initiation à l'algorithmique}
\author{Derrez Lorenzo}

\begin{document}
\maketitle

\section{Exercice n°1}
\subsection{Question 1 }
\subsubsection{En algorithme}
\begin{algorithm}
    \caption{Test}
    \begin{algorithmic}[1]
        \Procedure{Test}{$b$}
            \State $a \gets 1$
            \State $a \gets a+a$
            \State $T \gets$ allocArray$(10,3)$
            \State $T[3] \gets 2$
            \State $T[a] \gets T[a]+1$
            \For{$i = 1$ \textbf{to} $2$}
                \State $T[i+1] \gets T[i] + 1$
                \If{$T[i] = 4$}
                    \State \Return \textbf{true}
                \EndIf
            \EndFor
            \State \Return \textbf{true}
        \EndProcedure
    \end{algorithmic}
\end{algorithm}
\subsubsection{En C}
\begin{lstlisting}[language=C, caption={Exercice 1 question 1}]
bool test(int b) {
    int a=1;
    a=a+a;
    int* T=allocIntArray(10,3);
    T[3]=2;
    T[a]=T[a]+1;
    T[T[a]+T[a]]=T[a]+1;
    for(int i=1; i<=2; i++){
        T[i+1]=T[i]+1;
        if(T[i]==4)
            return 1;
    }
    return 1;
}
\end{lstlisting}

\subsection{Question 2}
Traduction $C \rightarrow$ Python I1 : \\
\begin{lstlisting}[language = Python, caption = {Exercice 1 question 2}]
def test(b:int) -> bool:
    a = 1
    a = a+a 
    T = newArray(10,3)
    T[3] = 2
    T[a] = T[a] + 1
    for i in range(1, 3): 
        T[i + 1] = T[i] + 1 
        if T[i] == 4: 
            return 1 
    return 1

\end{lstlisting}

\section{Exercice 2}
\subsection{Question 1}

\begin{itemize}
    \item Etant donné que le paramètre \textit{b} n'intervient pas dans le programme, la sortie du problème sera la même à chaque appel.
    \item La sortie en question est : 
        \subitem { [1, 4, 4, 3, 3, 3, 3, 3, 3, 3] }
        
\end{itemize}

\subsection{Question 2}

\begin{itemize}
    \item Le problème attendu est de chercher la première occurrence de 4 dans le tableau puis de le renvoyer. 
\end{itemize}

\section{Exercice 3}
\subsection{Question 1}

\begin{itemize}
    \item \underline{Jeu de test 1} :
        \\
        Entrée : x = 2; T = [1, 2, 4] 
        \\
        Sortie : (True, 1)
    \item \underline{Jeu de Test 2} :
        \\ 
        Entrée : x = 3; T = [1, 1, 2, 3, 1]
        \newline
        Sortie : (True, 3)
\end{itemize}

\subsection{Question 2}

\begin{itemize}
    \item \underline{Jeu de test 1} :
        \\
        Entrée : [1, 2, 4] \\
        Sortie : (False, $ \emptyset $)
    \item \underline{Jeu de test 2} :
        \\
        Entrée : [1, 1, 2, 3, 1]
        \\
        Sortie : (False, $ \varnothing $)
    \item \underline{Jeu de test 3} : 
        \\
        Entrée : [10, 3, 1]
        \\ Sortie : (True, $ \varnothing $)
\end{itemize}

\subsection{Question 3}

\begin{itemize}
    \item Le problème est de chercher si 10 est dans T et si oui noter son indice i dans un couple puis assigner 100 à l'indice i. 
    \\
    \begin{algorithm}
        \caption{Formalisme Logique de \textit{test2}}
        \begin{algorithmic}[1]
            \State Entrée : T un tableau d'entiers
            \State Sortie : $\varnothing$
            \State $(b,j) \leftarrow $ Trouver 10 dans T
            \If{b = True}
                \State $ T[j] \leftarrow $ 100
            \EndIf
        \end{algorithmic}
    \end{algorithm}
\end{itemize}

\subsection{Question 4}

\begin{algorithm}
    \caption{ : $1^e$ solution}
    \begin{algorithmic}[1]
        \State Entrée : \textit{x} entier, \textit{T} tableau d'entiers
        \State Sortie : un couple \textit{(b,i)} vérifiant la propriété de l'énoncé 
        \Function{Appartient}{$x, T$}
            \State (b,i) = isElement(x, T)
            \State \Return (b,i)
        \EndFunction
    \end{algorithmic}
\end{algorithm}

\begin{algorithm}
    \caption{ : $2^e$ solution}
    \begin{algorithmic}[1]
        \State Entrée : \textit{x} entier, \textit{T} tableau d'entiers
        \State Sortie : un couple \textit{(b,i)} vérifiant la propriété de l'énoncé
        \Function{Appartient}{$x, T$}
            \For{j allant de 0 à |T|-1}
                \If{ T[j] = x}
                    \State $ b \leftarrow $ True
                    \State \Return (b,i)
                \EndIf
            \EndFor
        \EndFunction
    \end{algorithmic} 
\end{algorithm}

\section{Exercice 4}

\subsection{Question 1}

\begin{algorithm}[H]
    \caption{yamMaxvalue \textit{(avec des fonctions auxiliaires)}}
    \begin{algorithmic}
        \State Entrée : \textit{T} Array[integer]
        \State Sortie : integer 
        \Function{yammaxvalue}{$T$}
            \State \Return $ max(yams(T), grande_suite(T), petite_suite(T), full(T), carré(T), brelan(T)) $
        \EndFunction
    \end{algorithmic}
\end{algorithm}

\subsection{Question 2}

\begin{algorithm}[H]
    \caption{\textbf{: Yams}}
    \begin{algorithmic}
        \State Entrée : \textit{T} tableau d'entiers
        \State Sortie : entier vérifiant : 
        \State \Return \textit{0} si pas vérifiée
        \State \Return \textit{50} si vérifiée
    \end{algorithmic}
\end{algorithm}

\begin{algorithm}[H]
    \caption{\textbf{: Grande Suite}}
    \begin{algorithmic}
        \State Entrée : \textit{T} tableau d'entiers
        \State Sortie : entier vérifiant : 
        \State \Return \textit{0} si pas vérifiée
        \State \Return \textit{40} si vérifiée
    \end{algorithmic}
\end{algorithm}

\begin{algorithm}[H]
    \caption{\textbf{: Petite suite}}
    \begin{algorithmic}
        \State Entrée : \textit{T} tableau d'entiers
        \State Sortie : entier vérifiant : 
        \State \Return \textit{0} si pas vérifiée
        \State \Return \textit{30} si vérifiée
    \end{algorithmic}
\end{algorithm}

\begin{algorithm}[H]
    \caption{\textbf{: Full}}
    \begin{algorithmic}
        \State Entrée : \textit{T} tableau d'entiers
        \State Sortie : entier vérifiant : 
        \State \Return \textit{0} si pas vérifiée
        \State \Return \textit{25} si vérifiée
    \end{algorithmic}
\end{algorithm}

\begin{algorithm}[H]
    \caption{\textbf{: Carré }}
    \begin{algorithmic}
        \State Entrée : \textit{T} tableau d'entiers
        \State Sortie : entier vérifiant : 
        \State \Return \textit{0} si pas vérifiée
        \State \Return \textit{Somme des quatres dés égaux} si vérifiée
    \end{algorithmic}
\end{algorithm}

\begin{algorithm}[H]
    \caption{\textbf{: Brelan}}
    \begin{algorithmic}
        \State Entrée : \textit{T} tableau d'entiers
        \State Sortie : entier vérifiant : 
        \State \Return \textit{0} si pas vérifiée
        \State \Return \textit{Somme des trois dés égaux} si vérifiée
    \end{algorithmic}
\end{algorithm}

\begin{algorithm}[H]
    \caption{\textbf{: Chance}}
    \begin{algorithmic}
        \State Entrée : \textit{T} tableau d'entiers
        \State Sortie : entier vérifiant : 
        \State \Return \textit{0} si pas vérifiée
        \State \Return \textit{Moitié$_{inf}$ de la somme des dés} si vérifiée
    \end{algorithmic}
\end{algorithm}

\section{Exercice 5}
\begin{itemize}
    \item Cela retourne le couple (100, 2).
\end{itemize}


\end{document}